% Source: https://artofproblemsolving.com/community/c3h2042850_x_and_ycomponents
% Stable identity: aops-topic-2042850-q01
Moving from \(X\) to \(Y\), the \(y\)-coordinate drops from \(4\) to \(-7.2\), a total change of \(11.2\). The point \(Z\) is where \(y=0\), which is reached after a drop of \(4\), i.e. after the fraction \(\dfrac{4}{11.2}=\dfrac{5}{14}\) of the way along the segment.

\begin{asy}
usepackage("amsmath");
size(5cm);
defaultpen(fontsize(10pt));
pair X=(3,4), Y=(9.6,-7.2), Z=(75/14,0);
draw((-0.5,0)--(11,0),Arrow(4)); draw((0,-8)--(0,5.5),Arrow(4));
draw(X--Y,blue+linewidth(0.9));
draw(X--(3,0),dashed); draw(Y--(9.6,0),dashed);
dot(X);dot(Y);dot(Z,red);
label("$X(3,4)$",X,NE); label("$Y(9.6,-7.2)$",Y,E); label("$Z$",Z,NE,red);
label("$4$",(3,2),W); label("$7.2$",(9.6,-3.6),E);
label("$x$",(11,0),E); label("$y$",(0,5.5),N);
label("\mbox{}",(min(currentpicture,user=true)+max(currentpicture,user=true))/2,black); // keeps the text after the figure black when LaTeX typesets the labels inline
\end{asy}

The \(x\)-coordinate increases by \(9.6-3=6.6\) along the whole segment, so
\[
x_Z=3+\dfrac{5}{14}\cdot 6.6=3+\dfrac{33}{14}=\boxed{\dfrac{75}{14}}.
\]
Since the distances along the segment are proportional to the changes in \(y\),
\[
XZ:ZY=4:7.2=\boxed{5:9}.
\]
(Check: the line has slope \(\dfrac{-11.2}{6.6}=-\dfrac{56}{33}\), and \(4-\dfrac{56}{33}\left(\dfrac{75}{14}-3\right)=4-\dfrac{56}{33}\cdot\dfrac{33}{14}=0\).)
